\scanpage{49}
\mspage{20}
diagramme d'inclusions \inserted{(croissantes)} d'ens tot. ordonnés
\[
(31)
\]
\diagram{\begin{tikzcd}[row sep=large,column sep=large]
&d''=d\cup d'&\\
d\arrow[ur,hook,no head]&&d'\arrow[ul,hook',no head]\\
&\delta=d\cap d'\arrow[ul,hook',no head]\arrow[ur,hook,no head]&
\end{tikzcd}}
tel que ce soit un diagramme cocart. dans la cat. des ens. ordonnés (p.ex $\delta$ segment initial de $d$, et $d'$ contenant l'élément final de $d$) alors le diagramme correspondant
\[
(32)
\]
\diagram{\begin{tikzcd}[row sep=large,column sep=large]
&\Sigma^{\unsure{<}}(d'')\arrow[dl]\arrow[dr]&\\
\Sigma^{<}(d)\arrow[dr]&&\Sigma^{<}(d')\arrow[dl]\\
&\Sigma^{\unsure{<}}(\delta)&
\end{tikzcd}}
est cartésien. \deleted{\unsure{(pour cela le cas est}}

Les conditions b c d du n\textsuperscript{o} précédent deviennent ici une seule (du type d)

\textcircled{b} Considérons un ens \deleted{\unsure{ens tot}} ordonné \inserted{de ``dim'' $\geq1$}, $d=(i_0,-i_n)$, et considérons le diagramme
\diagram{\begin{tikzcd}
\deleted{\unsure{\Sigma(n)}}\arrow[r,phantom,"\deleted{\longrightarrow}" description]&\deleted{\Sigma(0)}
\end{tikzcd}}
\[
(33)
\]
\diagram{\begin{tikzcd}[row sep=large,column sep=large]
&\Sigma^{<}(i_0,i_n)\arrow[dr,"\varphi_0"]&\\
\Sigma^{<}(d)\arrow[ur]\arrow[r]\arrow[ddr,bend right=15]&\Sigma^{<}(i_1,i_n)\arrow[r,"\varphi_1"]&\Sigma^{<}(i_n)\\
&\vdots&\\
&\Sigma^{<}(i_{n-1},i_n)\arrow[uur,"\varphi_{n-1}"']&
\end{tikzcd}}
\deleted{de \unsure{plongements} bordants.} d'immersions, où les $\varphi_0,\varphi_1,\ldots\varphi_n$ sont bordantes. On suppose
\margin{Plus précisément nous voulons que pour $x\in\Sigma^{\unsure{<}}(0)$, les $z_\alpha\in\Sigma^{<}(1)$ au dessus de $x$ (par but) sont en \uline{corr. biunivoque} avec les germes $B_\alpha$ des sous-ens. bord-ts, images locales des $\Sigma^{<}(1)$ en $z_\alpha$, que ces germes sont transversaux, et que les points \inserted{de $\Sigma^{\unsure{<}}(1)$} \deleted{\illegible{3}} \unsure{au dessus de $x$} \unsure{correspondants} \unsure{distincts} \unsure{déterminent} un drapeau \unsure{dans} \unsure{les $z_\alpha$}.}
\mspage{48}
